\documentclass[12pt]{article}
\usepackage[T1]{fontenc}
\usepackage[utf8]{inputenc}
\usepackage{lmodern}
\usepackage{textcomp}
\usepackage{enumitem}
\usepackage{geometry}
\geometry{margin=1in}
\begin{document}
\begin{center}
Answer Key - Health Sciences - K-12 Level Assessment 19 \
\small (Answers are ordered exactly as in the question paper.)
\end{center}

\section*{Multiple Choice}
\begin{enumerate}[label=\arabic*.]
\item \textbf{Answer:} C
\\ \textit{Options:} A) Brain growth is faster than at any other stage in life; B) Peak bone mass occurs before the growth spurt ends; C) Girls physiologically increase their percentage body fat; D) Type 2 obesity related diabetes remains a very rare problem
\\ \textit{Note:} During adolescence, girls undergo hormonal changes that promote the accumulation of body fat as part of normal physiological development, which is essential for reproductive health and overall growth.
\item \textbf{Answer:} A
\\ \textit{Options:} A) Visual performance; B) Motor performance; C) Behaviour disorders; D) Memory
\\ \textit{Note:} Supplementation with long-chain fatty acids EPA and DHA during pregnancy and infancy is linked to improved visual performance because these omega-3 fatty acids are crucial for the development of retinal and neural tissues in the growing brain.
\item \textbf{Answer:} D
\\ \textit{Options:} A) Brain; B) Kidney; C) Liver; D) Skeletal Muscle
\\ \textit{Note:} Skeletal muscle has the lowest energy expenditure per unit of body mass because, while it is metabolically active during physical activity, it uses less energy at rest compared to other organs like the brain or liver, which are consistently engaged in high-energy processes.
\item \textbf{Answer:} B
\\ \textit{Options:} A) The frequent use of antiseptic swipes will reduce allergies; B) Living on a working farm will help reducing the development of atopic diseases; C) Particular care must be taken during breast feeding in high risk infants; D) Living on a working farm will increase the development of atopic diseases
\\ \textit{Note:} The 'hygiene hypothesis' posits that increased exposure to microbes and allergens, such as those found on working farms, can help strengthen the immune system and reduce the risk of developing allergic conditions and atopic diseases.
\item \textbf{Answer:} C
\\ \textit{Options:} A) Is rapid in the first 10 weeks of foetal life; B) Is slow in infancy; C) In childhood, adolescence and young adulthood predicts risk of fracture in old age; D) Is particularly high after the menopause in women
\\ \textit{Note:} Bone mineral accretion during childhood, adolescence, and young adulthood is crucial because higher peak bone mass achieved during these stages is associated with a lower risk of fractures later in life due to improved bone density and strength.
\end{enumerate}

\section*{True / False}
\begin{enumerate}[label=\arabic*.]
\setcounter{enumi}{5}
\item \textbf{Answer:} False
\\ \textit{Options:} A) True; B) False
\\ \textit{Note:} Infants born to mothers who are vegan are not at increased risk of Vitamin C deficiency because a well-planned vegan diet can provide sufficient amounts of Vitamin C through fruits and vegetables that are rich in this essential nutrient.
\item \textbf{Answer:} True
\\ \textit{Options:} A) True; B) False
\\ \textit{Note:} Evaluating the safety of a food substance involves assessing not only its chemical composition and potential hazards but also the amount that individuals are likely to consume daily, as this determines the overall risk to human health.
\item \textbf{Answer:} True
\\ \textit{Options:} A) True; B) False
\\ \textit{Note:} The statement is true because regulatory agencies require that all pesticide chemicals, substances from food packaging, and color additives be rigorously tested for safety to protect public health before they can be approved for use in food products.
\item \textbf{Answer:} False
\\ \textit{Options:} A) True; B) False
\\ \textit{Note:} Muscle glycogen, not triacylglycerol, is primarily used as the energy source during short-term intense physical activities due to its quicker availability for rapid energy production.
\item \textbf{Answer:} True
\\ \textit{Options:} A) True; B) False
\\ \textit{Note:} The statement is true because hyperglycemia in type 2 diabetes results from a combination of the liver producing too much glucose, the pancreas not secreting enough insulin to regulate blood sugar levels, and muscle and fat cells becoming less responsive to insulin, leading to decreased glucose uptake.
\end{enumerate}

\section*{Long Form Response}
\begin{enumerate}[label=\arabic*.]
\setcounter{enumi}{10}
\item \textbf{Answer:} The Mediterranean diet, rich in fruits, vegetables, whole grains, legumes, nuts, and healthy fats like olive oil, is linked to preventing age-related cognitive decline. This diet emphasizes foods high in antioxidants and omega-3 fatty acids, which support brain health by reducing inflammation and oxidative stress. For example, studies have shown that individuals who closely follow the Mediterranean diet experience slower cognitive decline and a lower risk of Alzheimer's disease. Incorporating fish, which is a staple in this diet, provides essential omega-3 fatty acids that are crucial for maintaining brain function. Overall, the Mediterranean diet promotes a balanced intake of nutrients that can help safeguard cognitive abilities as we age.
\\ \textit{Note:} correct because it accurately describes how the Mediterranean diet's emphasis on nutrient-rich foods with antioxidants and omega-3 fatty acids can reduce inflammation and oxidative stress, which are linked to slower cognitive decline and a lower risk of Alzheimer's disease.
\item \textbf{Answer:} Plant sources play a significant role in providing essential amino acids for human nutrition. While all plant proteins contain all essential amino acids, some may have limitations in the quantity of specific amino acids. For example, legumes are often low in methionine, while grains may lack lysine. This means that while a diet based solely on one type of plant may not provide sufficient amounts of certain amino acids, combining different plant foods can help achieve a balanced amino acid profile. Therefore, it's important for individuals, especially those following a vegetarian or vegan diet, to consume a variety of plant proteins to ensure they receive all the essential amino acids needed for health.
\\ \textit{Note:} correct because it highlights that while plant sources can provide essential amino acids, certain plants may be deficient in specific amino acids, necessitating a combination of different plant foods to achieve a complete amino acid profile for optimal human nutrition.
\end{enumerate}

\vfill
\noindent\textit{End of Answer Key}
\end{document}